\chapter[\emj{🌿}~Recorridos de Árboles (Tree Traversals)]{\emj{🌿}~Recorridos de Árboles (Tree Traversals)}

\begin{dsaquees}

Visitar todas las habitaciones de una casa de varios pisos 🏠. Hay dos
grandes estilos:

\begin{itemize}
\item En profundidad (DFS): entras por una puerta y sigues bajando hasta
      el sótano antes de regresar. Según CUÁNDO anotes cada habitación
      (al entrar, en medio, o al salir) obtienes Preorder, Inorder o
      Postorder.
\item Por niveles (BFS): recorres piso por piso: primero toda la planta
      baja, luego todo el primer piso, y así. Eso es Level-order.
\end{itemize}

El "cuándo apuntas el nodo" cambia por completo el orden del resultado.

\end{dsaquees}

\begin{dsaotraforma}

Recorrer un árbol es visitar cada nodo exactamente una vez. Los cuatro
recorridos fundamentales:

\begin{itemize}
\item \textbf{Preorder} (Raíz → Izq → Der): útil para COPIAR/serializar
      un árbol.
\item \textbf{Inorder} (Izq → Raíz → Der): en un BST devuelve los
      valores ORDENADOS.
\item \textbf{Postorder} (Izq → Der → Raíz): útil para BORRAR el árbol o
      evaluar expresiones (procesa hijos antes que el padre).
\item \textbf{Level-order} (por niveles): BFS con Queue; base de muchos
      problemas de árboles en entrevistas.
\end{itemize}

Los tres DFS salen naturales con recursión, pero conviene dominar la
versión iterativa con Stack (evita stack overflow y la piden mucho).

\end{dsaotraforma}

\begin{dsadiagrama}
Árbol:
          [1]
        /     \
      [2]     [3]
     /   \       \
   [4]   [5]     [6]

Preorder  (R,I,D): 1  2  4  5  3  6
Inorder   (I,R,D): 4  2  5  1  3  6
Postorder (I,D,R): 4  5  2  6  3  1
Level-order      : 1  2  3  4  5  6   (por niveles, con Queue)
\end{dsadiagrama}

\begin{dsacomofunciona}
\begin{verbatim}
CUÁNDO se "apunta" el nodo cambia todo (mismo árbol):

          (1)
         /   \
       (2)   (3)
      /  \      \
    (4)  (5)    (6)

PREORDER  (apunta al ENTRAR, antes de bajar):
  visita 1 → baja izq → 2 → 4 → sube → 5 → sube → der → 3 → 6
  1 [2 (4)(5)] [3 (6)]           =  1 2 4 5 3 6

INORDER   (apunta ENTRE izq y der):
  4 (2) 5 ... 1 ... 3 6          =  4 2 5 1 3 6

POSTORDER (apunta al SALIR, después de los hijos):
  (4)(5) 2  (6) 3  1             =  4 5 2 6 3 1

LEVEL-ORDER (Queue, nivel por nivel):
  Queue: [1] → saca 1, mete 2,3 → [2,3] → saca 2, mete 4,5 ...
  nivel0: 1   nivel1: 2 3   nivel2: 4 5 6   =  1 2 3 4 5 6
\end{verbatim}
\end{dsacomofunciona}

\begin{lstlisting}
DFS RECURSIVO
- Preorder : procesa nodo, luego izq, luego der.
- Inorder  : izq, procesa nodo, der.
- Postorder: izq, der, procesa nodo.

LEVEL-ORDER (BFS)
- Queue con la raíz. Saca un nodo, procésalo, encola sus hijos.
- Para separar por niveles: procesa exactamente queue.size() nodos
  por iteración del bucle externo.

PSEUDOCÓDIGO (level-order por niveles)
──────────────────────────────────────────────────────────────

función levelOrder(root):
    SI root == NULL: return
    q = Queue con root
    MIENTRAS q no vacia:
        n = q.size()               // nodos de este nivel
        REPETIR n veces:
            nodo = q.dequeue()
            procesar(nodo)
            SI nodo.izq: q.enqueue(nodo.izq)
            SI nodo.der: q.enqueue(nodo.der)
        // aqui termina un nivel

CÓDIGO C++ (DFS recursivo)
──────────────────────────────────────────────────────────────

struct Node { int val; Node *left, *right; };

void preorder(Node* n)  { if(!n) return; cout<<n->val<<" ";
                          preorder(n->left);  preorder(n->right); }
void inorder(Node* n)   { if(!n) return; inorder(n->left);
                          cout<<n->val<<" "; inorder(n->right); }
void postorder(Node* n) { if(!n) return; postorder(n->left);
                          postorder(n->right); cout<<n->val<<" "; }

CÓDIGO C++ (inorder iterativo con Stack)
──────────────────────────────────────────────────────────────

#include <stack>
#include <vector>
using namespace std;

vector<int> inorderIter(Node* root) {
    vector<int> res; stack<Node*> st; Node* cur = root;
    while (cur || !st.empty()) {
        while (cur) { st.push(cur); cur = cur->left; } // baja a la izq
        cur = st.top(); st.pop();
        res.push_back(cur->val);                       // procesa
        cur = cur->right;                              // ve a la der
    }
    return res;
}

CÓDIGO C++ (level-order con Queue)
──────────────────────────────────────────────────────────────

#include <queue>
vector<vector<int>> levelOrder(Node* root) {
    vector<vector<int>> res;
    if (!root) return res;
    queue<Node*> q; q.push(root);
    while (!q.empty()) {
        int n = q.size();
        vector<int> nivel;
        for (int i = 0; i < n; i++) {
            Node* cur = q.front(); q.pop();
            nivel.push_back(cur->val);
            if (cur->left)  q.push(cur->left);
            if (cur->right) q.push(cur->right);
        }
        res.push_back(nivel);
    }
    return res;
}

COMPLEJIDAD (Big-O)
──────────────────────────────────────────────────────────────

  Recorrido     │ Tiempo  │ Espacio
  ──────────────┼─────────┼──────────────────────────
  Los 4         │ O(n)    │ O(h) DFS / O(w) BFS

* n = nodos, h = altura, w = anchura máxima del árbol.
\end{lstlisting}

\begin{dsaentrevistas}

✅ Inorder de un BST = ordenado: base para validar un BST o encontrar el
   k-ésimo menor. Si el inorder no sale creciente, NO es un BST válido.

💡 Level-order es la puerta a decenas de problemas: "Right Side View",
   "Zigzag Traversal", "conectar nodos del mismo nivel", "profundidad
   mínima". Casi siempre es BFS por niveles con \verb|queue.size()|.

⚠️ Recursión vs iterativo: la recursiva es más limpia, pero con árboles
   muy profundos (miles de niveles) puede desbordar el stack. Ten la
   versión iterativa lista.

🔑 Reconstrucción: con Preorder + Inorder (o Postorder + Inorder) puedes
   reconstruir un árbol binario único. Pregunta frecuente de Google.

\end{dsaentrevistas}

\begin{dsaresumen}

• Preorder (R,I,D): copiar/serializar.
• Inorder (I,R,D): en BST da orden creciente.
• Postorder (I,D,R): borrar árbol, evaluar expresiones.
• Level-order (BFS con Queue): recorridos por niveles.
• Domina la versión iterativa con Stack (evita stack overflow).
• Todos O(n) tiempo; espacio O(h) en DFS, O(ancho) en BFS.
════════════════════════════════════════════════════════════════

\end{dsaresumen}
